\section{到 NNLO QCD 的系数函数}\label{sec:coefs}

强耦合与夸克质量需要按
\begin{equation}
  \begin{split}
    g_0 = \left( \frac{\mu \, e^{\EulerGamma / 2}}{\sqrt{4 \pi}}
    \right)^\ep Z_g\,g\,,\qquad m_0 = Z_m\,m\,,
    \label{eq:msbar}
  \end{split}
\end{equation}
作常规 QCD 重正化，其中重正常数 $Z_g$ 与 $Z_m$ 写为
\begin{equation}
  \begin{split}
    Z_g &=
    1 - \gcoup{2}\frac{\beta_0}{2\ep} +
    \gcoup{4}\left(\frac{3\beta_0^2}{8\ep^2}
    -\frac{\beta_1}{4\ep}\right)   + \order{g^6}\,,
    \\
    Z_m &=
    1
    -\gcoup{2}\frac{\gamma_{m,0}}{2\ep}
    +\gcoup{4}\left[\frac{1}{\ep^2}\left(
      \frac{\gamma_{m,0}^2}{8} + \frac{\beta_0\gamma_{m,0}}{4}\right)
      -\frac{\gamma_{m,1}}{4\ep}\right] + \order{g^6}\,,
    \label{eq:rgconsts}
  \end{split}
\end{equation}
且
\begin{equation}
  \begin{split}
      \beta_0 &= \frac{11}{3} C_A - \frac{4}{3} \TF\,, \qquad
      \beta_1 = \frac{34}{3} C_A^2 - \left( 4 C_F + \frac{20}{3} C_A
      \right) \TF \,,\\
      \gamma_{m,0} &= 6 C_F\,,\qquad
      \gamma_{m,1} = \frac{97}{3} C_A C_F + 3 C_F^2
      - \frac{20}{3} C_F \TF \,.
    \label{eq:anombg}
  \end{split}
\end{equation}
$C_F$ 与 $C_A$ 分别是规范群基本表示与伴随表示的二次 Casimir 本征值。此外，
$\TF=Tn_F$，其中 $n_F$ 是夸克味的数目，$T$ 是基本表示中迹的归一化。对
SU($N_c$)，有 $C_F=(N_c^2-1)/(2N_c)$、$C_A=N_c$、$T=1/2$。

另外，流动夸克场也需要按
\begin{equation}
  \begin{split}
    \chi_{f,R} = \sqrt{Z_\chi}\,\chi_{f}\,,
  \end{split}
\end{equation}
重正化，这导致 \eqn{eq:tops} 中算符 $\tcalo_{3,4}$ 定义里的因子 $Z_\chi$。
$Z_\chi$ 不同于常规 QCD 的夸克场重正常化。\nlo\ 的 \msbar\ 结果已在
\citere{Luscher:2013cpa} 中求出。为确定我们计算所需的 \nnlo\ 重正常数，
可以利用 $c_3(t)$ 在重正化后必须有限这一事实。把 \msbar\ 表达式写成
\begin{equation}
  \begin{split}
    Z^{-1}_\chi &=
    1
      -\gcoup{2}\frac{\gamma_{\chi,0}}{2\ep}
      +\gcoup{4}\left[\frac{1}{\ep^2}\left(
        \frac{\gamma_{\chi,0}^2}{8} +
        \frac{\beta_0\gamma_{\chi,0}}{4}\right)
        -\frac{\gamma_{\chi,1}}{4\ep}\right] + \order{g^6}\,,
    \label{eq:zchi}
  \end{split}
\end{equation}
我们发现（脚注：注意由于领头阶 $c_4=0$，该系数只需要 \nlo\ 重正化。）
\begin{equation}
  \begin{split}
  \gamma_{\chi,0} &= 6C_F\,,\\
  \gamma_{\chi, 1} &=
  C_A C_F \left( \frac{223}{3} - 16 \ln 2 \right) - C_F^2 \left( 3 + 16
  \ln 2 \right) - \frac{44}{3} C_F \TF\,.
  \end{split}
\end{equation}
这使得我们能在 \msbar\ 方案下求出到 \nnlo\ QCD 的能动量张量系数（脚注：为
方便读者，我们还随论文以电子附件文件的形式提供了
$c_1,\ldots,c_4$ 的表达式。）：
\begin{equation}
  \begin{split}
    c_1(t) =& \, \frac{1}{g^2} \Bigg\{ 1 + \gcoup{2}
    \left[ - \frac{7}{3}  C_A + \frac{3}{2}
      \TF -  \beta _0 \, L (\mu, t) \right] \\
	& \quad + \gcoup{4} \Bigg[ - \beta_1 \, L
      (\mu, t) + C_A^2 \left( - \frac{14482}{405} - \frac{16546}{135}
      \ln 2 + \frac{1187}{10} \ln 3 \right) \\
      & \qquad + C_A \TF \Bigg( \frac{59}{9}
      \text{Li}_2\left(\frac{1}{4}\right) + \frac{10873}{810} +
      \frac{73}{54} \pi^2 - \frac{2773}{135} \ln 2 + \frac{302}{45} \ln
      3 \Bigg) \\
	& \qquad + C_F \TF \Bigg( - \frac{256}{9}
      \text{Li}_2\left(\frac{1}{4}\right) +\frac{2587}{108} -
      \frac{7}{9} \pi^2 - \frac{106}{9} \ln 2 - \frac{161}{18} \ln 3
      \Bigg) \Bigg] \\
    & \quad + \order{g^6} \Bigg\} \, ,
    \label{eq:c1}
  \end{split}
\end{equation}
\begin{equation}
	\begin{split}
          c_2(t) =& \, \frac{1}{4 g^2} \Bigg\{ - 1 +  \gcoup{2}
          \left[ \frac{25}{6}  C_A - 3  \TF + \beta _0 \,
      L (\mu, t) \right] \\
	& \quad +  \gcoup{4} \Bigg[ \beta _1 \, L (\mu, t) +
      C_A^2 \left(\frac{56713}{1620} - \frac{1187}{10} \ln 3 +
      \frac{16546}{135} \ln 2 \right) \\
	& \qquad + C_A  \TF \Bigg( - \frac{59}{9}
      \text{Li}_2\left(\frac{1}{4}\right) - \frac{6071}{405} -
      \frac{73}{54} \pi^2 + \frac{2287}{135} \ln 2 - \frac{361}{90} \ln
      3 \Bigg) \\
	& \qquad + C_F \TF \Bigg( \frac{220}{9}
      \text{Li}_2\left(\frac{1}{4}\right) - \frac{1757}{54} +
      \frac{10}{9} \pi ^2 - \frac{164}{9} \ln 2 + \frac{247}{9} \ln 3
      \Bigg) \Bigg] \\
	& \quad  + \order{g^6} \Bigg\} \, ,
    \label{eq:c2}
  \end{split}
\end{equation}
\begin{equation}
	\begin{split}
	  c_3 (t) =& \, \frac{1}{4} \Bigg\{1 + \gcoup{2}
          \left( \frac{3}{2} C_F +
          \frac{\gamma_{\chi, 0}}{2} L (\mu, t) \right) \\
			& \quad + \gcoup{4}
          \Bigg[ \frac{\gamma_{\chi, 0}}{4} \left(\beta _0+
            \frac{\gamma_{\chi, 0}}{2} \right) \Big( L^2 (\mu, t) + L
            (\mu, t) \Big) + \frac{\gamma_{\chi, 1}}{2} L (\mu, t) \\
			& \qquad + C_F^2 \Bigg( - \frac{137}{9}
            \text{Li}_2\left(\frac{1}{4}\right) - \frac{559}{216} +
            \frac{103}{108} \pi^2 - \frac{1736}{27} \ln 2 +
            \frac{122}{3} \ln 3 - 4 \ln^2 2 \Bigg) \\
			&  \qquad + C_F  \TF \Bigg( -
            \frac{136}{9}  \text{Li}_2\left(\frac{1}{4}\right) -
            \frac{3377}{810} - \frac{7}{9} \pi^2 + \frac{1232}{135} \ln
            2 - \frac{136}{15} \ln 3 \Bigg) \\
			& \qquad + C_A C_F \Bigg( - \frac{365}{9}
            \text{Li}_2 \left(\frac{1}{4}\right) + \frac{261829}{3240} +
            \frac{77}{108} \pi^2 + \frac{5788}{45} \ln 2 \\
			& \qquad \quad - \frac{2102}{15} \ln 3 - 4 \ln^2
            2 \Bigg) \Bigg] \\
			& \quad + \order{g^6} \Bigg\} \, ,
    \label{eq:c3}
	\end{split}
\end{equation}
\begin{equation}
  \begin{split}
    c_4 (t) =& \, \frac{C_F}{2} \, \Bigg\{ \gcoup{2} +
    \gcoup{4} \Bigg[ \left( \beta _0
      + \frac{\gamma_{\chi, 0}}{2} \right) L (\mu, t) \\
      & \qquad + C_F \Bigg( - \frac{161}{18}
      \text{Li}_2\left(\frac{1}{4}\right) - \frac{41}{54} -
      \frac{55}{108} \pi^2 - \frac{1105}{27} \ln 2 +
      \frac{101}{6} \ln 3 \Bigg) \\
      & \qquad + \TF \Bigg( \frac{25}{9}
      \text{Li}_2\left(\frac{1}{4}\right) -
      \frac{20573}{1620} + \frac{5}{18} \pi^2 +
      \frac{6559}{135} \ln 2 - \frac{679}{30} \ln 3 \Bigg)
      \\
      & \qquad + C_A \Bigg( \frac{257}{36}
      \text{Li}_2\left(\frac{1}{4}\right) - \frac{137}{405}
      + \frac{11}{216} \pi^2 - \frac{419}{90} \ln 2 +
      \frac{1157}{60} \ln 3 \Bigg) \Bigg] \\
    & \quad + \order{g^6} \Bigg\} \, ,
    \label{eq:c4}
  \end{split}
\end{equation}
其中引入了参数（脚注：该参数的动机来自流时间积分中出现的典型因子
$(8\pi t)^{\ep}$\,\cite{Luscher:2010iy} 与 \msbar\ 方案中重正化标度的惯用定义
（见 \eqn{eq:msbar}）之乘积：$(8\pi t)^\ep
(\mu^2e^{\EulerGamma}/(4\pi))^\ep = 1 +
  \ep\,L(\mu,t) + \order{\ep^2}$。）
\begin{equation}
	L (\mu, t) \equiv \ln \left( 2 \mu^2 t \right) + \EulerGamma\,,
        \label{eq:lmut}
\end{equation}
其中 Euler--Mascheroni 常数为 $\EulerGamma=0.57721\ldots$。在 \nlo\ 层面，
这些结果与 \citere{Suzuki:2013gza,Makino:2014taa} 完全一致。我们在常规 QCD
的一般 $R_\xi$ 规范中进行计算；规范参数依赖在最终结果中消去，这是另一项
受欢迎的检验。\eqn{eq:flow} 中的规范参数 $\kappa$ 取为 1。

虽然能动量张量 $T_{\mu\nu}(x)$ 与重正化方案无关，但对算符
$\tcalo_{i,\mu\nu}(t,x)$ 和系数函数 $c_i(t)$ 却未必如此。由于
$\tcalo_{1,\mu\nu}$ 与 $\tcalo_{2,\mu\nu}$ 无需算符重正化，其矩阵元以及相应
系数函数确实与重正化方案无关。另一方面，若按 \msbar\ 方案使用 \eqn{eq:zchi}
的夸克场重正常数 $Z_\chi$，则对 $i\in\{3,4\}$，$\tcalo_{i,\mu\nu}$ 的矩阵元
和系数函数 $c_i(t)$ 会显式依赖重正化标度 $\mu$。

不过，这种方案依赖可以通过引入所谓``带圈''（ringed）夸克场来避免，如
~\citere{Makino:2014taa} 所建议。这相当于把 \eqn{eq:tops} 中的 $Z_\chi$
替换为
\begin{equation}
  \begin{split}
    \mathring{Z}_\chi(t) = \frac{-2 N_c \, n_F}{(4 \pi t)^2
      \langle\bar{\chi}_f(t,x)\overleftrightarrow
      {\slashed D}\chi_f(t,x)\rangle}\,.
    \label{eq:zf}
  \end{split}
\end{equation}
目前 $\mathring{Z}_\chi(t)$ 只知到 \nlo\ QCD。不过，其显式的 $\mu$ 依赖项可以
根据 $c_i(t)$ 对 $i\in \{3,4\}$ 必须有限且与 $\mu$ 无关的要求加以重构。
这样我们得到它与 \eqn{eq:zchi} 的 \msbar\ 夸克场重正常数之比：
\begin{equation}
	\begin{split}
	  \zeta_\chi&\equiv \frac{\mathring{Z}_\chi}{Z_\chi} = 1 + \gcoup{2}
          \left( \frac{\gamma_{\chi, 0}}{2} L(\mu ,t) - 3 C_F \ln 3 - 4
          C_F \ln 2 \right) \\ &
                \qquad + \gcoup{4} \Bigg\{
                  \frac{\gamma_{\chi, 0}}{4} \left( \beta_0 +
                  \frac{\gamma_{\chi, 0}}{2} \right) L^2 (\mu, t) +
                  \Big[ \frac{\gamma_{\chi, 1}}{2} -
                  \frac{\gamma_{\chi, 0}}{2} \left( \beta_0 +
                  \frac{\gamma_{\chi, 0}}{2} \right) \ln 3 \\
                & \qquad \qquad -
                  \frac{2}{3}
                  \gamma_{\chi, 0} \left(
                  \beta_0 + \frac{\gamma_{\chi, 0}}{2}\right) \ln 2
                  \Big] L(\mu,t)
                  + C_2 \Bigg\} \\ &\qquad + \order{g^6} \,.
                \label{eq:zfchi}
	\end{split}
\end{equation}
常数 $C_2$ 无法用这种方式确定，而需要对 \eqn{eq:zf} 分母中出现的两点函数做
专门的单圈之外的三圈计算。本文不详细描述该计算；它将与对我们框架更完整的描述
一起在后续出版物\,\cite{Artz:2019bpr}中给出。此处我们只引用保留三位有效数字的
数值结果——考虑到下文讨论的理论不确定性，这已绰绰有余（脚注：注意
$\mathring{Z}_\chi$ 的计算也为 $Z_\chi$ 提供了独立检验。）：
\begin{equation}
  C_2 = -23.8 \, C_A C_F + 30.4 \, C_F^2 - 3.92 \, C_F T_F\,.
\end{equation}
把 Eqs.\,(\ref{eq:c3}) 和 (\ref{eq:c4}) 中的 $c_3(t)$、
$c_4(t)$ 乘以此比值，也使这些系数在形式上与 $\mu$ 无关，即
\begin{equation}
  \begin{split}
    \mu\dderiv{}{\mu}\{c_1,c_2,\ringc_3,\ringc_4\}=0\,,\quad
    \text{其中}\quad \ringc_i\equiv \zeta_\chi^{-1} c_i\,.
  \end{split}
\end{equation}
与任何微扰计算一样，这种 $\mu$ 无关性只在 $g$ 的更高阶精度内成立。剩余
$\mu$ 依赖的减小因此常用作对所研究可观测量的微扰展开的定性检验。于是我们
研究四个系数在把 $c_3$、$c_4$ 除以 \eqn{eq:zfchi} 定义的比值 $\zeta_\chi$
之后的 $\mu$ 依赖。我们固定一个特征流时间 $t$，让重正化标度 $\mu$ 围绕中心值
$\mu_0$ 变化；$\mu_0$ 由条件 $L(\mu_0,t)=0$（见 \eqn{eq:lmut}）定义，即
\begin{equation}
  \begin{split}
    \mu_0 = \frac{e^{-\EulerGamma/2}}{\sqrt{2t}}\,.
    \label{eq:mu0}
  \end{split}
\end{equation}
图\,\ref{fig:mu3gev} 与 \ref{fig:mu130gev} 展示了 $c_1$、$c_2$、$\ringc_3$、
$\ringc_4$ 的领头阶（\lo）、\nlo\ 与 \nnlo\ 近似作为重正化标度的函数在两个
不同流时间 $t$（分别对应 $\mu_0=3$\,GeV 与 $\mu_0=130$\,GeV）下的行为。
前者取 $n_F=3$，后者取 $n_F=5$。我们用 $\alpha^{(n_F=5)}_s(M_Z)=0.118$
求出耦合的输入值
$g^{(n_F=3)}(3\,\text{GeV})=1.77$ 与
$g^{(n_F=5)}(130\,\text{GeV})=1.19$。强耦合常数 $g(\mu)$ 的 $\mu$ 变化通过
借助 \texttt{RunDec}\,\cite{Chetyrkin:2000yt,Herren:2017osy} 数值求解相应的
重正化群方程确定：\lo、\nlo、\nnlo\ 曲线分别在单圈、双圈、三圈水平求解。
在 \fig{fig:mu3gev} 中，$t$ 的选取使 \eqn{eq:mu0} 的中心标度为
$\mu_0=3$\,GeV。在该中心标度处，\nnlo\ 修正使系数 $c_1$ 与 $c_2$ 的模相对
NLO 分别增大 10\% 与 13\%。这在把 $\mu/\mu_0$ 于 1/2 到 2 之间变分所估计的、
因缺失更高阶效应导致的 \nlo\ 不确定性的两倍以内——那里 $c_1$ 为 $7.3$\%、
$c_2$ 为 $8.0\%$。因此我们有信心认为：用同样方式估计的 \nnlo\ 不确定性相当
可靠——它对 $c_1$ 为 $5.7$\%，对 $c_2$ 为 $7.2$\%。注意这些数字的主要贡献
来自 $\mu$ 向下的变分，那里 $g(\mu)$ 开始对非微扰区敏感。$c_1$ 与 $c_2$ 在
较大 $\mu$ 处的行为似乎提示：这一不确定性估计实际上可能过于保守。


\begin{figure}
  \begin{center}
    \begin{tabular}{cc}
      \includegraphics[width=.45\textwidth,bb=20 0 300 220]{images/c1_3.pdf} &
      \includegraphics[width=.45\textwidth,bb=20 0 300 220]{images/c2_3.pdf} \\
      \includegraphics[width=.45\textwidth,bb=20 0 300 220]{images/c3_3.pdf} &
      \includegraphics[width=.45\textwidth,bb=20 0 300 220]{images/c4_3.pdf}
    \end{tabular}
  \end{center}
	\caption{系数 $c_1$、$c_2$、$\ringc_3=\zeta^{-1}_\chi c_3$、
          $\ringc_4=\zeta^{-1}_\chi c_4$（定义见
          \eqs{eq:c1}--\noeqn{eq:c4}，$\zeta_\chi$ 见
          \eqn{eq:zfchi}）的重正化标度依赖。黑色点线、蓝色虚线与红色实线分别对应于
          在 $c_1$、$c_2$ 中保留至 $(g^2)^{n-1}$ 阶、在 $\ringc_3$ 与
          $\ringc_4$ 中保留至 $(g^2)^{n}$ 阶的项，$n=0,1,2$。中心标度取为
          $\mu_0=3$\,GeV，对应 $t = 3.12
          \cdot 10^{-2} /\text{GeV}^2$，见 \eqn{eq:mu0}。味的数目取
          $n_F=3$。}
	\label{fig:mu3gev}
\end{figure}


\begin{figure}
\begin{center}
    \begin{tabular}{cc}
      \includegraphics[width=.45\textwidth,bb=20 0 300 220]{images/c1_130.pdf} &
      \includegraphics[width=.45\textwidth,bb=20 0 300 220]{images/c2_130.pdf} \\
      \includegraphics[width=.45\textwidth,bb=20 0 300 220]{images/c3_130.pdf} &
      \includegraphics[width=.45\textwidth,bb=20 0 300 220]{images/c4_130.pdf}
    \end{tabular}
  \end{center}
	\caption{同 \fig{fig:mu3gev}，但 $\mu_0=130$\,GeV（或
          $t=1.66\cdot 10^{-5} /\text{GeV}^2$），$n_F=5$。}
	\label{fig:mu130gev}
\end{figure}


与胶子型系数 $c_1$、$c_2$ 相反，费米子算符的系数 $c_3$、$c_4$ 只有从
\nlo\ 项开始才表现出剩余标度依赖。因此可以预期这两项在 \nnlo\ 有更强的
$\mu$ 依赖。尽管如此，对 $\ringc_3$，因标度变分导致的理论不确定性估计仍从
$9.8$\% 降到 $8.1$\%。\nnlo\ 效应带来的结果增大在 $\mu=\mu_0$ 处相对 \nlo\
结果为 $8.3$\%。

另一方面，$\ringc_4$ 在 $\mu_0=3$\,GeV 处的行为不那么令人满意。此时
\nnlo\ 效应使 \nlo\ 结果增加一倍以上，而且从 \nlo\ 到 \nnlo，因标度变分给出的
不确定性估计实际上从 47\% \textit{升}到 71\%。但请注意 \lo\ 时 $c_4=0$，
这意味着该系数在数值上是次要的。

正如预期，对 $\mu_0=130$\,GeV，所有系数的微扰行为显著改善，参见
\fig{fig:mu130gev}。对 $c_1$、$c_2$ 与 $\ringc_3$，标度不确定性在 \nlo\
已处于亚百分点水平；到 \nnlo，三种情形都小于 $0.2$\%。\nnlo\ 修正相对
\nlo\ 结果的效应约为：$c_1$、$c_2$ 各约 $2$\%，$\ringc_3$ 为 $0.8$\%。对
$\ringc_4$ 情况也明显好转：\nnlo\ 项使 \nlo\ 结果增加 38\%，不确定性从
\nlo\ 的 10\% 降到 \nnlo\ 的 $5.8$\%。

同样值得指出的是，\eqn{eq:mu0} 所定义的中心标度 $\mu_0$ 的选取看起来得到了
逐阶行为的支持。在几乎所有情形，NLO 与 NNLO 修正在 $\mu=\mu_0$ 处
\textit{都}相对较小。同时，NNLO 修正相对 NLO 结果总是小于 NLO 修正相对
LO 结果。唯一例外是 $\mu_0=3$\,GeV 处的 $\ringc_4$，但那里也没有哪个
$\mu$ 的选取显得格外突出。

综上，我们的结论是：NNLO 项显著改善了 Wilson 系数的微扰精度。
